\section{Arquitectura y direccionamiento}

\begin{table}[H]
  \centering
  \begin{tabularx}{\linewidth}{l l l Y}
    \toprule
    \textbf{Nodo} & \textbf{Sistema} & \textbf{IP interna} & \textbf{Función} \\
    \midrule
    \texttt{pve-lab5} & Proxmox VE & 10.50.5.1 & Hipervisor y administración de VM \\
    \texttt{lab5-cliente} & openSUSE Leap & 10.50.5.10 & Navegador, SSH, FileZilla y resolución local \\
    \texttt{lab5-web} & Ubuntu Server & 10.50.5.20 & Apache, PHP, WordPress y cliente MySQL \\
    \texttt{lab5-db} & Ubuntu Server & 10.50.5.30 & MySQL y base \texttt{wordpress} \\
    \bottomrule
  \end{tabularx}
  \caption{Direccionamiento propuesto para la red interna.}
\end{table}

La red \texttt{10.50.5.0/24} permanece pendiente de validación contra la red física. Proxmox utilizará \texttt{vmbr1} como puente interno. Cada VM podrá disponer de una segunda interfaz para actualizaciones, pero solo una interfaz tendrá ruta predeterminada.

\begin{center}
\fbox{\parbox{0.9\linewidth}{\centering
\textbf{openSUSE 10.50.5.10}\\
\vspace{0.3em}
$\downarrow$ HTTPS 443 y SSH/SFTP 22222\\
\vspace{0.3em}
\textbf{WordPress 10.50.5.20}\\
\vspace{0.3em}
$\downarrow$ MySQL 13306\\
\vspace{0.3em}
\textbf{MySQL 10.50.5.30}}}
\end{center}

\pendiente{Confirmar nombres de interfaces, subred y rutas después de crear las VM.}
